\section{一定角の回転}\label{節：一定角の回転}

単位円上の点を$z_n=e^{i\theta_n}$と書き,\,一定角$\phi$ずつ進めると
\[
 z_{n+1}=e^{i\phi}z_n,\qquad
 \theta_n=\theta_1+(n-1)\phi
\]
となる.
横座標$a_n=\cos\theta_n$だけを記録すると,\,加法公式から
\[
 \cos(\theta_n+2\phi)+\cos\theta_n
 =2\cos\phi\cos(\theta_n+\phi)
\]
なので
\[
 a_{n+2}=2\cos\phi\,a_{n+1}-a_n.
\]
一つの横座標だけでは点の上下を区別できないが,\,隣り合う二つの記録があれば次を更新できる.
特性方程式の根は$e^{i\phi},e^{-i\phi}$であり,\sectionref{節：隣接3項間漸化式と回転}につながる.
$0<\phi<\pi$では二根が非実数となる.

一方,\,傾き$a_n=\tan\theta_n$を記録すれば
\[
 a_{n+1}=\frac{a_n+\tan\phi}{1-a_n\tan\phi}
\]
という一次分数型になる.
ただし$\tan\phi$と各項の正接が定義され,\,更新の分母が非零の範囲で考える.
同じ角度の動きでも,\,\textbf{横座標と傾きでは更新式の形が変わる}.
